\documentclass{article}

% NeurIPS 2026 style, preprint option: no line numbers, authors shown.
\usepackage[preprint]{neurips_2026}

\usepackage[utf8]{inputenc}
\usepackage[T1]{fontenc}
\usepackage{hyperref}
\usepackage{url}
\usepackage{booktabs}
\usepackage{amsfonts}
\usepackage{nicefrac}
\usepackage{microtype}
\usepackage{xcolor}
\usepackage{graphicx}
\usepackage{adjustbox}
\usepackage{array}
\usepackage{enumitem}
\usepackage[most]{tcolorbox}
\usepackage[capitalise,noabbrev]{cleveref}

% Macros written by the analysis pipeline. Fallbacks so that the report compiles before it has run.
\IfFileExists{generated/tables/macros.tex}{\input{generated/tables/macros}}{}
\providecommand{\DataGenerated}{(no data yet)}
\providecommand{\NRunsReal}{0}
\providecommand{\NRunsFailed}{0}
\providecommand{\NRunsSynthetic}{0}

% Question box at the top of each analysis section.
\newtcolorbox{question}{enhanced, breakable, colback=white, colframe=black!35,
  boxrule=0.4pt, leftrule=2pt, arc=0pt, left=6pt, right=6pt, top=3pt, bottom=3pt,
  fontupper=\small}
% Analysis still to write (shown in color so that it can't be mistaken for results).
\newcommand{\placeholder}[1]{{\color{orange!80!black}\textit{[To write: #1]}}}
% Figure that needs data not collected yet (all of it, or part of it, synthetic).
% Shown only when the figure was drawn with synthetic data (flags written by the pipeline).
\IfFileExists{generated/figures/synthetic.tex}{\input{generated/figures/synthetic.tex}}{}
\makeatletter
\newcommand{\syntheticnote}[1]{\expandafter\ifx\csname syn@#1\endcsname\relax\else
  \if1\csname syn@#1\endcsname{\color{red!70!black}\textbf{Synthetic data}: placeholder values from a
  toy model, not measurements; hollow markers.}\ \fi\fi}
\makeatother
% Generated figure, or a box saying its data is missing (the figure file is deleted when it has no data).
\newcommand{\fig}[2][\linewidth]{\IfFileExists{generated/figures/#2.pdf}%
  {\includegraphics[width=#1]{generated/figures/#2}}%
  {\fbox{\parbox{0.9\linewidth}{\centering\small\vspace{1em}No data yet for \texttt{\detokenize{#2}}.\vspace{1em}}}}}
% Value from DESIGN.md that is an estimate, not a measurement.
\newcommand{\est}{\textit{(est.)}}

\title{Large-scale LLM pre-training\\ on Trillium's B200 partition}

\author{%
  Milabench team \\
  Mila -- Quebec AI Institute \\
}

\begin{document}

\maketitle

\begin{abstract}
  We evaluate the B200 partition of Trillium (SciNet) for large-scale LLM pre-training: 52 nodes of
  8 B200 GPUs (416 GPUs), with as much inter-node bandwidth per FLOP as Rorqual but nodes twice as large.
  Two workloads are used: Qwen3.8-27B (dense, TorchTitan) and GLM-5 (744B mixture of experts, Megatron-LM).
  We report node and fabric health, interference between jobs running side by side, the choice of
  parallel layout for each workload, scaling to the full partition, the gain of FP8 and MXFP8, and
  operational aspects (checkpointing, stability, startup, energy).
  \placeholder{key numbers once the window is analyzed.}
\end{abstract}

\begin{tcolorbox}[colback=red!4, colframe=red!50!black, boxrule=0.4pt, arc=0pt, fontupper=\small]
  \textbf{Draft generated \DataGenerated.} \NRunsReal{} runs measured on Trillium
  (\NRunsFailed{} failed), and \NRunsSynthetic{} synthetic runs standing in for experiments not run yet.
  Every figure with synthetic data says so and draws it with hollow markers.
\end{tcolorbox}

\section{Introduction}
\label{sec:intro}

Trillium's B200 partition gives exclusive access to 52 nodes of 8 B200 GPUs for a few days
\est{}. We use the window to measure how well the partition trains large language models, and which
configurations make the best use of it. \Cref{tab:machine} lists what is known about the machine before
access.

\begin{table}[h]
  \caption{Trillium's B200 partition, as known before access.}
  \label{tab:machine}
  \centering\small
  \begin{adjustbox}{max width=\linewidth}
  \begin{tabular}{lll}
    \toprule
    & Known & Source \\
    \midrule
    Nodes & 52 (\texttt{trig0101}--\texttt{trig0152}), exclusive access during the window & SciNet talk 2026-10-21; \texttt{/etc/hosts} \\
    Per node & 8 B200 180\,GB on NVLink, 2 $\times$ 48-core Intel CPUs, 2\,TB RAM, liquid cooled & SciNet talk \\
    GPUs & 416 in total & \\
    Network & InfiniBand, 800\,Gb/s per node, ``non-blocking''; number of HCAs unknown & user; SciNet \\
    Switch layout & unknown; \texttt{/etc/hosts} puts the nodes in 8 subnets of 7, 7, 7, 7, 4, 7, 7, 6 nodes & \texttt{/etc/hosts} \\
    Shared with others & distributed filesystem and interconnect (CPU and H100 nodes stay in production) & user \\
    Software & \texttt{module load cuda/12.9}, \texttt{--gpus-per-node=b200:8} & SciNet talk \\
    Matmul ($8192^2$) & B200: BF16 1{,}668, FP8 3{,}041\,TFLOP/s; H100: BF16 785\,TFLOP/s & SciNet (their measurement) \\
    Window & a few days \est{}, date unknown & user \\
    \bottomrule
  \end{tabular}
  \end{adjustbox}
\end{table}

\paragraph{What sets the partition apart.}
\begin{enumerate}[leftmargin=*, itemsep=1pt]
  \item \textbf{Network per FLOP.} Each GPU has 100\,Gb/s of inter-node bandwidth for 1{,}668\,TFLOP/s
        of BF16 matmul (SciNet's measurement): about 0.060\,Gb/s per TFLOP/s. This is close to Rorqual
        (50\,Gb/s for 785\,TFLOP/s, 0.064), where the network limited training, but the NVLink domain is
        twice as large (8 GPUs instead of 4). A standard DGX B200 has 4 times more network per node.
        Configurations that keep most traffic inside a node should do best.
  \item \textbf{FP8 and MXFP8.} Low precision doubles the matmul rate (3{,}041\,TFLOP/s in FP8, SciNet's
        measurement) but not the network. What is network-bound in BF16 becomes more so.
  \item \textbf{Memory.} 180\,GB per GPU, 73\,TB in total. Qwen3.8-27B (about 430\,GB of training state
        \est{}) fits on one node; GLM-5 fits on about 13 nodes \est{}, so up to 4 copies run at once.
        Activation checkpointing can be reduced.
  \item \textbf{Unknown topology, shared fabric.} The switch layout is not known, and the filesystem
        and interconnect carry other users' traffic. Placement enters the design only if the pairwise
        measurements show a structure (\cref{sec:network}); a reference run repeated through the window
        separates our configurations from background load (\cref{sec:interference}).
\end{enumerate}

\paragraph{Questions.} Each section answers one question:
(\cref{sec:network}) are all GPUs and links alike, and does the fabric have a structure?
(\cref{sec:interference}) can jobs run side by side without slowing each other?
(\cref{sec:layouts}) which parallel layout should each workload use?
(\cref{sec:scaling}) how far does each workload scale on the partition?
(\cref{sec:precision}) what do FP8 and MXFP8 buy, and do they train the same?
(\cref{sec:operations}) what do checkpoints, startup and energy cost, and is throughput stable?

\paragraph{Declared deviations} from the reference training recipes, valid for every result unless noted:
\placeholder{list once the B200 configurations are fixed (optimizer, multi-token prediction, data).}

\section{Methodology}
\label{sec:method}

\paragraph{Workloads.} \Cref{tab:workloads} summarizes the two training workloads. Both train every
parameter (memory allows it; no frozen experts in GLM-5).

\begin{table}[h]
  \caption{Workloads. The Trillium additions were tested on other clusters' H100 with truncated models, not yet on B200.}
  \label{tab:workloads}
  \centering\small
  \begin{adjustbox}{max width=\linewidth}
  \begin{tabular}{lll}
    \toprule
    & Qwen3.8-27B (dense) & GLM-5 (744B MoE, 256 experts, top-8) \\
    \midrule
    Framework & TorchTitan & Megatron-LM \\
    Parallelism today & FSDP, HSDP, TP (at most 4: 4 KV heads) & EP, PP, DP; TP $=1$ and CP $=1$ \\
    Added for Trillium & PP, MXFP8, gradient-reduction dtype, checkpointing, loss per step
           & 8 GPUs per node, checkpoint save / load, FP8 / MXFP8, B200 layouts, loss per step \\
    Precision recipes & BF16, MXFP8 & BF16, FP8 tensorwise, FP8 blockwise, MXFP8 \\
    Sequence & $3\times4{,}096$ per GPU (12{,}288 tokens) by default & 8{,}192, micro-batch 1 (the only one possible) \\
    Global batch & grows with the GPUs (weak scaling) & weak: 256 / 512 / 768 / 1{,}024 at 13 / 26 / 39 / 52 nodes; \\
                 & & strong: 768 at all four; small-batch point: 64 at 13 nodes \\
    Recompute & selective by default; full, none & full by default; selective (MoE activation and norms) \\
    \bottomrule
  \end{tabular}
  \end{adjustbox}
\end{table}

The pinned TorchTitan version removed its Float8 tensorwise and rowwise recipes (2026-10-01), so Qwen
compares BF16 with MXFP8 only. On B200, GLM-5's FP8 blockwise recipe runs as MXFP8 matrix multiplications.
GLM-5 stays at micro-batch 1: its plugin's attention takes THD input only, so the global batch is chosen in micro-batches.
With $P$ stages and $m$ micro-batches per data-parallel rank, the 1F1B pipeline bubble is $(P-1)/(m+P-1)$ \est{} (the
formula, to compare with the runs). Weak scaling keeps 32 micro-batches per data-parallel rank, a bubble of 27\,\%; strong scaling
keeps 768 sequences in total: 96, 48, 32 and 24 micro-batches, bubbles of 11, 20, 27 and 33\,\%; the small batch of 64 has 8 micro-batches, 60\,\%.

\paragraph{Network arithmetic for Qwen \est{}.} Any data-parallel layout (FSDP, HSDP, TP + FSDP) moves
about 8 bytes per parameter per node per step: BF16 parameter gathers in the forward and backward
passes, and an FP32 gradient reduction. That is about 215\,GB, or at least 2.2\,s per step at 800\,Gb/s,
against about 3\,s of compute at 12{,}288 tokens per GPU and 50\,\% of 1{,}668\,TFLOP/s. The levers that
change this ratio are the tokens per GPU, a BF16 gradient reduction (about 6 bytes per parameter), and
pipeline parallelism across nodes (volume divided by the number of stages). HSDP instead of FSDP changes
when the traffic happens, not its volume.

\paragraph{Layouts.} Qwen layouts are named after their parallel dimensions: \texttt{fsdp} shards over
all GPUs; \texttt{hsdp8} and \texttt{hsdp16} shard within 1 or 2 nodes and replicate across;
\texttt{tp2-fsdp} and \texttt{tp4-fsdp} add tensor parallelism of 2 or 4; \texttt{pp2-fsdp} and
\texttt{pp4-fsdp} add 2 or 4 pipeline stages. GLM-5 layouts are written \texttt{pp}$P$\texttt{-ep}$E$:
$P$ pipeline stages and expert groups of $E$ GPUs. Megatron carves the expert groups out of each stage's
data-parallel ranks: on $N$ nodes, $8N$ must be divisible by $P$, and $8N/P$ by $E$. A model copy is
therefore $P\times E$ GPUs. With 13 stages (the 78 layers split 7, then 11 stages of 6, then 5: the first
stage holds the dense layers, the last the loss) and EP8, one copy uses $13\times8=104$ GPUs, 13 nodes,
so 4 copies fill 52 nodes (\texttt{pp13-ep8}). Its memory: Rorqual's full GLM-5 run peaked at
66.6\,GiB on each of 200 GPUs, about 13\,TiB; 13 B200 nodes hold about 17\,TiB \est{}.

\paragraph{Throughput.} Qwen runs skip 30 warm-up steps (compilation, allocator growth), then measure for
5 minutes. GLM-5 runs train 20 iterations and leave out the first. The official throughput is the median
of the per-step global token rates of the measured steps, reported per GPU (tokens/s per GPU).

\paragraph{MFU} $= \text{FLOPs per token} \times \text{tokens/s per GPU} / \text{peak}$, with the dense
BF16 peak of the B200 specification, $2.25\times10^{15}$\,FLOP/s per GPU, for every precision recipe. For
context, SciNet measured 1{,}668\,TFLOP/s for an $8192^2$ BF16 matmul, 74\,\% of that peak. MFU is
derived from a measured rate and an analytic FLOP count.

\paragraph{Telemetry and health.} Every node of every job is sampled once a second by the node monitor
(GPU clocks, throttle reasons, power, temperature and duty cycle, XID and ECC errors, NVLink, InfiniBand,
CPU clocks). Every node is checked before use (\cref{sec:network}); failing nodes are excluded for the rest
of the window, and a node is used whole or not at all.

\paragraph{Reference runs and repetitions.} A reference run (Qwen, 4 fixed nodes, FSDP, 12{,}288 tokens
per GPU, 10 minutes) is repeated about every 3 hours through the window, to tell changes in our
configurations from changes in the shared fabric and filesystem (\cref{sec:interference}). Layout
experiments are repeated twice; checkpoint experiments three times at different times of day. Four GLM-5 runs are not
run a second time and serve two experiments: \texttt{pp13-ep8} at 13 and 52 nodes (A-G, also weak scaling), at 39 nodes (weak, also
strong scaling) and at 26 nodes (weak, also the layout reference at 26 nodes).

\paragraph{Measured, estimated and synthetic values.}
\begin{itemize}[leftmargin=*, itemsep=1pt]
  \item \emph{Measured}: filled markers, solid lines and bars.
  \item \emph{Estimated} (derived from a model or from other runs, e.g.\ the compute part of the step in
        \cref{fig:step_decomposition}): ``(est.)'' in the label or legend.
  \item \emph{Synthetic} (placeholder for an experiment not run yet): hollow markers, dashed lines or
        hollow bars, a ``synthetic'' watermark on the figure, and a note in its caption. Numbers in the
        text that come from synthetic data are shown as \TriSyn{1.0}.
\end{itemize}
Estimates never replace measurements; both are kept in the data. Synthetic values are removed as the
experiments they stand in for are run.

\paragraph{Data.} Raw results are written to \texttt{\$SCRATCH/trillium/} on Trillium and copied
read-only to the analysis machine; every figure, table and number here is regenerated from them by one
command. Formats: \texttt{scripts/analysis/FORMAT.md}, with the Trillium additions in
\texttt{scripts/trillium/FORMAT.md} (layout and precision fields, loss per step, checkpoint events,
collective sweeps). \Cref{tab:ledger} lists every run, including failures. Only runs on Trillium enter the analysis.

\section{Health and network}
\label{sec:network}

\begin{question}
  \textbf{Are all GPUs and links alike, and does the fabric have a structure?} Do all 416 B200 deliver the
  same matmul rate? Does every node reach its NVLink and InfiniBand bandwidth? Do some node pairs
  communicate faster than others? Does collective bandwidth hold up to 52 nodes?
\end{question}

This group runs first: it excludes faulty nodes, and it decides whether node placement is a factor of
the later experiments.

\paragraph{GPUs (D1).} A per-node health check measures, on every GPU, the BF16 and FP8 matmul rate, and
on every node the NVLink all-reduce bandwidth, the loopback bandwidth of each HCA and the GPU--HCA
affinity. \Cref{fig:health_gpus} compares every GPU with the median.

\begin{figure}[h]
  \centering\fig{health_gpus}
  \caption{Matmul rate of every GPU in BF16 and FP8, by node (strip), with the median and a line 5\,\%
    below it. \syntheticnote{health_gpus}}
  \label{fig:health_gpus}
\end{figure}

\placeholder{median rates against SciNet's numbers; GPUs below the $-5\,\%$ line and what was done with them.}

\paragraph{Links (D1).} \Cref{fig:health_links} shows the intra-node and per-HCA bandwidths.

\begin{figure}[h]
  \centering\fig{health_links}
  \caption{Per-node NVLink all-reduce bus bandwidth (left) and per-HCA bandwidth (right).
    \syntheticnote{health_links}}
  \label{fig:health_links}
\end{figure}

\placeholder{number of HCAs per node and their rate; outlier nodes or links.}

\paragraph{Pairs (D2).} Every pair of nodes is measured: 26 disjoint pairs at a time, 51 rounds, all
1{,}326 pairs. If the matrix shows groups of nodes (latency or bandwidth clusters, for example along the
8 subnets of \texttt{/etc/hosts}), a placement experiment is added (Qwen and GLM-5 at 13 nodes inside one
group vs.\ spread over groups). Otherwise placement is dropped from the design.

\begin{figure}[h]
  \centering\fig{pair_matrix}
  \caption{Pairwise bandwidth (and latency) between the 52 nodes, ordered by subnet.
    \syntheticnote{pair_matrix}}
  \label{fig:pair_matrix}
\end{figure}

\placeholder{whether the matrix has a structure, and the decision on placement.}

\paragraph{Collectives at scale (D3, D4).} NCCL all-reduce, all-gather, reduce-scatter and all-to-all
with 1\,GiB messages at 1, 2, 4, 8, 13, 26 and 52 nodes, and a message-size sweep at 2 and 52 nodes. The
bisection test (D4) splits the nodes into two random halves of 26 that exchange data all-to-all, for
3 random cuts. A non-blocking fabric keeps the per-node bandwidth constant as nodes are added.

\begin{figure}[h]
  \centering\fig{collectives_scale}
  \caption{Bus bandwidth vs.\ number of nodes, one line per collective, with the 800\,Gb/s line rate of a
    node; bisection points (D4). \syntheticnote{collectives_scale}}
  \label{fig:collectives_scale}
\end{figure}

\placeholder{bandwidth at 52 nodes as a share of the line rate; bisection bandwidth; whether the fabric
behaves as non-blocking.}

\section{Running jobs side by side}
\label{sec:interference}

\begin{question}
  \textbf{Can jobs run side by side without slowing each other, and does the background load move our
  numbers?} If four 13-node jobs run as fast together as alone, the later experiments can be packed in
  parallel. The fabric and filesystem are shared with other partitions, so the same run may also change
  speed during the window.
\end{question}

\paragraph{Interference (E1).} The Qwen best guess (HSDP within a node, 24k tokens per GPU) runs on 4
disjoint random sets of 13 nodes: each set alone, then the 4 together. This is repeated twice with new
random sets. If the slowdown together is below 2\,\%, runs of up to 26 nodes are packed in parallel for
the layout, scaling, precision and operations experiments. The median change together vs.\ alone is
\TriInterferencePct\,\%.

\begin{figure}[h]
  \centering\fig[0.72\linewidth]{interference}
  \caption{Tokens/s per GPU of each 13-node set (2 draws of 4 sets), alone (circles) and with the 3
    others running (squares). \syntheticnote{interference}}
  \label{fig:interference}
\end{figure}

\placeholder{whether any set slows down beyond the run-to-run spread; decision on packing.}

\paragraph{Reference drift (R).} The reference run (\cref{sec:method}) is repeated about every 3 hours.
Its spread bounds the differences between configurations that can be attributed to the configurations
themselves: the fastest reference run is \TriDriftPct\,\% faster than the slowest.

\begin{figure}[h]
  \centering\fig[0.72\linewidth]{reference_drift}
  \caption{Throughput of the reference run (Qwen, 4 fixed nodes, FSDP, 12{,}288 tokens per GPU) vs.\ time
    since the first reference run. Band: median $\pm1\,\%$; the linear trend over the window is written
    on top. \syntheticnote{reference_drift}}
  \label{fig:reference_drift}
\end{figure}

\placeholder{correlation with time of day or with other users' load.}

\section{Parallel layouts}
\label{sec:layouts}

\begin{question}
  \textbf{Which parallel layout should each workload use on this partition?} For Qwen3.8-27B: which layout,
  at which batch per GPU, and which lever (batch, gradient-reduction dtype, pipeline parallelism,
  activation checkpointing) buys the most? How much of the step is spent in exposed communication? For
  GLM-5: should expert parallelism stay inside a node or cross nodes?
\end{question}

\paragraph{Qwen3.8-27B (A-Q).} The factors and levels are in \cref{tab:aq}. The full cross of layouts and
tokens per GPU runs at 13 nodes; the 3 best layouts also run at 2 and 52 nodes, and the layouts without
pipeline parallelism at 1 node. About 45 runs, each repeated twice, about 10 minutes each \est{}.

\begin{table}[h]
  \caption{Factors of the Qwen layout experiment (A-Q).}
  \label{tab:aq}
  \centering\small
  \begin{adjustbox}{max width=\linewidth}
  \begin{tabular}{ll}
    \toprule
    Factor & Levels \\
    \midrule
    Layout & \texttt{fsdp} / \texttt{hsdp8} / \texttt{hsdp16} / \texttt{tp2-fsdp} / \texttt{tp8-fsdp} /
             \texttt{pp2-fsdp} / \texttt{pp4-fsdp} \\
    Tokens per GPU & 12k ($4096\times3$) / 24k ($4096\times6$) / 48k ($4096\times12$) \\
    Gradient reduction & FP32 / BF16 (on the 2 best layouts) \\
    Activation checkpointing & selective / full / none (1 and 13 nodes; none only where it fits) \\
    Nodes & 13 (full cross); 2 and 52 for the 3 best; 1 for the layouts without PP \\
    \bottomrule
  \end{tabular}
  \end{adjustbox}
\end{table}

\begin{figure}[h]
  \centering\fig{qwen_layouts}
  \caption{Qwen3.8-27B throughput (tokens/s per GPU, or MFU) by layout (rows) and tokens per GPU
    (columns), at 2, 13 and 52 nodes (panels). \syntheticnote{qwen_layouts}}
  \label{fig:qwen_layouts}
\end{figure}

\placeholder{best layout and batch at each scale; whether keeping the shards in a node (\texttt{hsdp8},
\texttt{tp8-fsdp}) wins as expected.}

\begin{figure}[h]
  \centering\fig{qwen_levers}
  \caption{Throughput change of each lever (batch, gradient-reduction dtype, pipeline parallelism,
    activation checkpointing) relative to \texttt{fsdp} at 12k tokens per GPU, at 13 nodes.
    \syntheticnote{qwen_levers}}
  \label{fig:qwen_levers}
\end{figure}

\placeholder{ranking of the levers; comparison with the network arithmetic of \cref{sec:method}.}

\paragraph{Where the step time goes.} \Cref{fig:step_decomposition} splits the step time into compute,
taken from 1-node runs at the same batch per GPU, and exposed communication, the remainder of the
measured step time. With 8 GPUs on NVLink, the 1-node runs train the full model, so the compute part is
measured, not scaled from a truncated model.

\begin{figure}[h]
  \centering\fig[0.85\linewidth]{step_decomposition}
  \caption{Step time of Qwen3.8-27B split into compute (1-node step time) and exposed communication, by
    layout and number of nodes. \syntheticnote{step_decomposition}}
  \label{fig:step_decomposition}
\end{figure}

\placeholder{share of exposed communication by layout at 13 and 52 nodes.}

\paragraph{GLM-5 (A-G).} The global batch is 64 sequences per model copy. The layouts differ in where
the expert groups lie (\cref{tab:ag}). Recomputation (full vs.\ selective) and micro-batch size (1 vs.\
2) are varied on the best layout. Each runs at 13 nodes (1 copy) and 52 nodes (4 copies). About 14 runs,
each repeated twice, about 25 minutes each \est{}, including about 15 minutes of startup at scale \est{}.

\begin{table}[h]
  \caption{GLM-5 layouts (A-G).}
  \label{tab:ag}
  \centering\small
  \begin{tabular}{llll}
    \toprule
    Layout & Expert group & Nodes per copy & Copies on 52 nodes \\
    \midrule
    \texttt{pp13-ep8}  & in a node (NVLink) & 13 & 4 \\
    \texttt{pp13-ep16} & across 2 nodes & 13 & 4 \\
    \texttt{pp13-ep32} & across 4 nodes & 13 & 4 \\
    \texttt{pp26-ep8}  & in a node & 26 & 2 \\
    \texttt{pp6-ep16} (13 layers per stage) & across 2 nodes & 12 & 4 (48 nodes), if it fits \\
    \bottomrule
  \end{tabular}
\end{table}

\begin{figure}[h]
  \centering\fig{glm_layouts}
  \caption{GLM-5 throughput (tokens/s per GPU) and all-to-all time by layout, at 13 and 52 nodes.
    \syntheticnote{glm_layouts}}
  \label{fig:glm_layouts}
\end{figure}

\placeholder{cost of expert groups crossing nodes; best layout, recomputation and micro-batch.}

\section{Scaling}
\label{sec:scaling}

\begin{question}
  \textbf{How far does each workload scale on the partition?} How does Qwen3.8-27B's throughput per GPU
  change from 8 to 416 GPUs, with the best layout and with plain FSDP? Do GLM-5 model copies scale from
  1 to 4?
\end{question}

All scaling runs use the best layouts found in \cref{sec:layouts}. In weak scaling the work per GPU is
fixed, so perfect scaling is a flat line of throughput per GPU. Unlike on 4-GPU nodes, a 1-node run
trains the full Qwen model, so it serves as the 100\,\% efficiency reference.

\paragraph{Qwen3.8-27B (B-Q).} Weak scaling at 1, 2, 4, 8, 13, 26 and 52 nodes with fixed tokens per GPU,
for the best layout and the \texttt{fsdp} baseline (B-Q-weak). Strong scaling keeps the global batch of
the full partition ($52\times8\times24$k tokens) at 13, 26 and 52 nodes (B-Q-strong).

\begin{figure}[h]
  \centering\fig{qwen_scaling}
  \caption{Qwen3.8-27B tokens/s per GPU and scaling efficiency vs.\ number of GPUs (8--416), best layout
    vs.\ \texttt{fsdp}; strong scaling in a separate panel. \syntheticnote{qwen_scaling}}
  \label{fig:qwen_scaling}
\end{figure}

\placeholder{efficiency at 52 nodes for both layouts; where the curve bends; strong-scaling efficiency.}

\paragraph{GLM-5 (B-G).} 1, 2, 3 and 4 model copies (13, 26, 39 and 52 nodes), with 64 sequences per
copy (weak scaling) and 64 sequences in total (strong scaling).

\begin{figure}[h]
  \centering\fig{glm_scaling}
  \caption{GLM-5 tokens/s per GPU vs.\ number of model copies (1--4), weak vs.\ strong batch.
    \syntheticnote{glm_scaling}}
  \label{fig:glm_scaling}
\end{figure}

\placeholder{whether copies interfere through the shared fabric; strong-scaling efficiency.}

\section{Precision}
\label{sec:precision}

\begin{question}
  \textbf{What do FP8 and MXFP8 buy, and do they train the same?} Low precision doubles the matmul rate
  but not the network: does the gain measured on one node survive at \TriNNodes{} nodes? Does the loss follow BF16, within its noise?
\end{question}

\paragraph{Speed (C-speed).} Recipes: BF16 and MXFP8 for Qwen3.8-27B (the pinned TorchTitan has no other
FP8 recipe); BF16, FP8 tensorwise, FP8 blockwise and MXFP8 for GLM-5. Each runs on the best layout, at 1
node (Qwen) or 13 nodes (GLM-5), and at the largest scale (\TriNNodes{} nodes, GLM-5 \TriGlmTopNodes{}). For Qwen, \TriQwenFpBestRecipe{} speeds up training by
$\times$\TriQwenFpSpeedupSmall{} on 1 node and $\times$\TriQwenFpSpeedupFull{} on \TriNNodes{} nodes. For GLM-5, the
best recipe, \TriGlmFpBestRecipe, gives $\times$\TriGlmFpSpeedupSmall{} on 13 nodes and
$\times$\TriGlmFpSpeedupFull{} on \TriGlmTopNodes{} nodes.

\begin{figure}[h]
  \centering\fig{precision_speed}
  \caption{Speedup over BF16 (ratio of tokens/s) by recipe, at small scale (1 node for Qwen3.8-27B, 13
    for GLM-5) and on the whole machine. \syntheticnote{precision_speed}}
  \label{fig:precision_speed}
\end{figure}

\placeholder{how much of the matmul gain survives in a training step, and how much the network takes
back at scale.}

\paragraph{Loss (C-loss).} The same recipes train for 500 steps at 13 nodes, from the same seed, on the
same FineWeb data, with a realistic learning rate. The difference to a BF16 run at each step measures
the effect of the recipe, if it stands out of the noise: BF16 runs twice more, once with the same seed and once with another,
and those repeats give the band of \cref{fig:precision_loss}. GLM-5's learning rate starts at 0 and warms up over 50 steps: steps
1 and 2 see the same weights, and an unstable phase follows, so its first steps say little about a recipe.

\begin{figure}[h]
  \centering\fig{precision_loss}
  \caption{Loss difference to one BF16 run vs.\ step, by recipe, at 13 nodes: Qwen3.8-27B (left), GLM-5 (right).
    Thick lines: 25-step mean; thin lines: per step; legend: mean over the last 100 steps; gray band: the same difference for the BF16 repeats
    (same seed, another seed), the noise. Dotted line: end of GLM-5's learning-rate warm-up. The vertical axis fits the settled part.
    \syntheticnote{precision_loss}}
  \label{fig:precision_loss}
\end{figure}

\placeholder{loss difference at step 500 against the band; whether a same-seed repeat is closer than another seed.}

\section{Operations}
\label{sec:operations}

\begin{question}
  \textbf{What does it cost to run production training here?} How long do checkpoints take to save and
  load, and how long is training blocked? Is throughput stable over an hour, and are some GPUs or nodes
  slower? How long until training starts? How many tokens does a kWh buy?
\end{question}

\paragraph{Checkpoints (F1).} Save and load of Qwen3.8-27B (about 0.38\,TB \est{}) at 1, 13 and 52 nodes,
and of GLM-5, weights only (about 1.5\,TB \est{}) and full (about 10.4\,TB \est{}), at 13 and 52 nodes;
synchronous vs.\ asynchronous saves (time training is blocked vs.\ total write time); 3 repetitions at
different times of day. The nodes have no local disk, so an asynchronous save can only stage in host
memory before writing to the shared VAST filesystem. That filesystem serves the rest of Trillium, so the
sizes escalate with the administrators' agreement (Qwen 1 node, Qwen 52 nodes, GLM-5 weights, GLM-5 full),
and only one GLM-5 checkpoint is on disk at a time, deleted once its load is timed.

At the largest scale, a synchronous save takes \TriCkptQwenSave\,s for Qwen3.8-27B and
\TriCkptGlmFullSave\,s for the full GLM-5 state; an asynchronous save of GLM-5 blocks training for
\TriCkptGlmFullBlocked\,s. The highest write throughput is \TriCkptWriteGBs\,GB/s.

\begin{figure}[h]
  \centering\fig{checkpoint}
  \caption{Checkpoint time vs.\ number of nodes for Qwen3.8-27B (full state), GLM-5 weights only and GLM-5
    full state: synchronous save, asynchronous save (total and time training is blocked), and load.
    Right: filesystem throughput of saves and loads. \syntheticnote{checkpoint}}
  \label{fig:checkpoint}
\end{figure}

\placeholder{variation with time of day; write throughput against the filesystem's rated bandwidth.}

\paragraph{Stability (F2).} The best Qwen layout trains for at least one hour on 52 nodes. GPU clocks and
power from the node telemetry share the time axis with throughput, so a drop can be traced to a GPU.

\begin{figure}[h]
  \centering\fig{stability}
  \caption{Long run (52 nodes), vs.\ time: tokens/s per GPU per step with a 15-step median (top); SM
    clock (middle) and power (bottom) of the GPUs: gray band, range over all 416 GPUs; line, the median;
    the GPU with the lowest value is named. \syntheticnote{stability}}
  \label{fig:stability}
\end{figure}

\placeholder{throughput drift over the hour; throttling events and their cause.}

\paragraph{Stragglers (F2).} With data parallelism every rank waits for the slowest one, so per-rank step
times are equal and cannot point to a slow GPU. The slow GPU is the one that waits least: its activity is
highest. A node whose CPU clock falls also holds back every rank on it.

\begin{figure}[h]
  \centering\fig{stragglers}
  \caption{Long run, from the node telemetry. Top: activity of each GPU (highest 10-minute median), by
    node. Bottom: per node, the lowest 1-minute median of the CPU busy-core clock relative to the first
    minute. Nodes grouped by subnet; outliers are named. \syntheticnote{stragglers}}
  \label{fig:stragglers}
\end{figure}

\placeholder{slowest GPUs and nodes, and their effect on throughput.}

\paragraph{Startup (F3).} Time from process start to the first step, taken from every run (no extra
runs). At the largest scale it is \TriStartupQwen\,min for Qwen3.8-27B and \TriStartupGlm\,min for
GLM-5.

\begin{figure}[h]
  \centering\fig[0.72\linewidth]{startup}
  \caption{Time from process start to the first step vs.\ number of nodes, for both workloads.
    \syntheticnote{startup}}
  \label{fig:startup}
\end{figure}

\placeholder{which phase (environment, initialization, compilation, first step) grows with scale.}

\paragraph{Energy (F4).} Tokens per kWh, from the GPU power sampled by the node telemetry in every run (no
extra runs). Only GPU power is counted: CPUs, memory, network and cooling are not. The most efficient
Qwen layout at 13 nodes reaches \TriEnergyBest{} million tokens per kWh.

\begin{figure}[h]
  \centering\fig{energy}
  \caption{Million tokens per kWh (GPU power only): Qwen3.8-27B by layout at 13 nodes and 24k tokens per
    GPU (left); by workload and precision on the whole machine (middle); Qwen3.8-27B weak scaling vs.\
    nodes (right). \syntheticnote{energy}}
  \label{fig:energy}
\end{figure}

\placeholder{energy per token of the best configurations; effect of precision and scale.}

\section{Conclusions and recommendations}
\label{sec:conclusion}

\placeholder{answers to the six questions, in order; recommended layout, batch and precision for each
workload; recommendations for SciNet (network, checkpoints, health checks) and for future evaluations.}


\appendix
\section{Run plan}
\label{app:plan}

The experiment groups run in the order of \cref{tab:plan}. Every time and count in this appendix is an
estimate made before access. If the 5 nodes reserved for SciNet's inference service are not
released, 47 nodes are available and the 52-node runs shrink accordingly (39 nodes for 3 GLM-5 copies).

\begin{table}[h]
  \caption{Experiments. Times are estimates.}
  \label{tab:plan}
  \centering\small
  \begin{adjustbox}{max width=\linewidth}
  \begin{tabular}{lp{0.62\linewidth}ll}
    \toprule
    ID & What & Scale & Time \est{} \\
    \midrule
    D1 & Per-node health check: per-GPU matmul BF16 / FP8, NVLink all-reduce, per-HCA loopback bandwidth, GPU--HCA affinity & every node & 30\,min \\
    D2 & Pairwise latency and bandwidth: 26 disjoint pairs per round, 51 rounds (all 1{,}326 pairs) & 52 nodes & 15\,min \\
    D3 & NCCL all-reduce, all-gather, reduce-scatter, all-to-all, 1\,GiB, at 1--52 nodes; size sweep at 2 and 52 nodes & 1--52 & 45\,min \\
    D4 & Bisection: two random halves of 26 all-to-all with each other, 3 cuts & 52 & 15\,min \\
    P  & Only if D2 shows groups: Qwen and GLM-5 at 13 nodes inside one group vs.\ spread & 13 & -- \\
    \midrule
    E1 & Qwen (HSDP-8, 24k tokens/GPU) on 4 disjoint sets of 13 nodes, alone then together; 2 repetitions & 52 & 1.5\,h \\
    R  & Reference: Qwen, 4 fixed nodes, FSDP, 12{,}288 tokens/GPU, 10\,min, every $\sim$3\,h & 4 & 10\,min each \\
    \midrule
    A-Q & Qwen layouts $\times$ tokens per GPU $\times$ reduction dtype $\times$ checkpointing (\cref{tab:aq}); $\sim$45 runs $\times$ 2 $\times$ $\sim$10\,min & 1--52 & see \cref{tab:budget} \\
    A-G & GLM-5 layouts (\cref{tab:ag}), 256 sequences per copy; recompute, micro-batch, 64 sequences per copy; $\sim$14 runs $\times$ 2 $\times$ $\sim$25\,min & 12--52 & see \cref{tab:budget} \\
    \midrule
    B-Q-weak & Qwen, best layout and \texttt{fsdp}, fixed tokens per GPU & 1--52 & \\
    B-Q-strong & Qwen, fixed global batch ($52\times8\times24$k tokens) & 13, 26, 52 & \\
    B-G & GLM-5 \texttt{pp13-ep8}, 1--4 copies, 256 sequences per copy (weak) and in total (strong) & 13--52 & \\
    \midrule
    C-speed & Qwen: BF16 / MXFP8; GLM-5: BF16 / FP8 tensorwise / FP8 blockwise / MXFP8; best layout & 1 or 13, and 52 & \\
    C-loss & Same recipes, 500 steps, same seed and FineWeb data, realistic learning rate & 13 & \\
    \midrule
    F1 & Checkpoint save / load, sync vs.\ async (staging in host memory: no local disk), 3 repetitions; escalation agreed with the administrators & 1--52 & \\
    F2 & Long run: best Qwen layout, $\geq$1\,h & 52 & \\
    F3 & Startup, from every run & -- & -- \\
    F4 & Energy, from every run & -- & -- \\
    \bottomrule
  \end{tabular}
  \end{adjustbox}
\end{table}

\begin{table}[h]
  \caption{Budget \est{}. ``Packed'' applies if E1 shows less than 2\,\% slowdown, so that runs of up to
    26 nodes run in parallel.}
  \label{tab:budget}
  \centering\small
  \begin{tabular}{lrr}
    \toprule
    Group & Serial (est.) & Packed (est.) \\
    \midrule
    Bring-up on B200 (install, 1-node smoke tests, FP8 / MXFP8 smoke tests) & 4\,h & \\
    D & 2\,h & \\
    E + R & 2.5\,h & \\
    A-Q & 15\,h & 6\,h \\
    A-G & 12\,h & 6\,h \\
    B & 8\,h & 5\,h \\
    C & 6\,h & 4\,h \\
    F & 6\,h & \\
    \midrule
    Total & $\sim$55\,h & $\sim$35\,h \\
    \bottomrule
  \end{tabular}
\end{table}

\paragraph{Priorities if the window is short.} Tier 0: bring-up, D1. Tier 1: A-Q at 13 nodes, B-Q-weak,
A-G at 13 and 52 nodes, B-G. Tier 2: D2--D4, E1, C, F1. Tier 3: F2, B-Q-strong, the rest of A.

\section{Run ledger}
\label{app:ledger}

\begin{table}[h]
  \caption{Run ledger: every attempt on Trillium, including failures.}
  \label{tab:ledger}
  \centering
  \begin{adjustbox}{max width=\linewidth}
    \IfFileExists{generated/tables/ledger.tex}{\input{generated/tables/ledger}}{\textit{No runs yet.}}
  \end{adjustbox}
\end{table}


\end{document}
